% Paragraph 1: What we did
We proposed a novel precision medicine approach utilizing multi-utility value functions for survival outcomes in the presence of competing risks  to estimate optimal individualized treatment regimes. Our main contributions are: 1) estimating optimal ITRs based on OS and priority cause cumulative incidence simultaneously, 2) prioritizing maximum OS over minimum priority cause risk, 3) introducing a multi-utility value function, and 4) developing random survival and cumulative forests that output individualized CIF curves for single stage treatment allocation. Under standard assumptions and mild regularity conditions, we established uniform consistency of the value function of our optimal ITR estimator, as well as for the proposed survival and cumulative incidence random forests.

% Paragraph 2: Difference in Literature – This is probably too long.
Our work builds on the work of \cite{cho_multi-stage_2023}, which is designed for finite decision points for survival outcomes and do not consider cause-specific events.
Additionally, our work extends utility functions for multiple outcomes \citep{Butler2018,jiang_precision_2021,zitovsky2023} by presenting a multi-utility approach. Our work is closest to \cite{jiang_precision_2021}, but they used multiple utility functions to construct a composite of equally weighted outcomes, while we implement it directly, incorporating hierarchical structure through a two-step procedure.
With respect to the existing CR PM methods \citep{friedenreich_physical_2016, yavuz_estimating_2018,chen_cumulative_2020,morzywolek_timing_2022}, our work most closely relates to \cite{zhou_restricted_2021}, who minimized side effects in a CR setting by estimating the optimal ITR at a specific time point, with side effects constrained below a set threshold. While our method is similar, we incorporate a pre-specified tolerance to determine whether the CR is applied. Additionally, \cite{zhou_restricted_2021} only considered linear decision rules whereas our approach uses flexible forest-based ensemble methods to capture complex data patterns. We also compared to \cite{he_optimal_2021}, who’s approach is doubly robust, while our nonparametric method entirely avoids the risk of model misspecification by not relying on parametric assumptions. \cite{he_optimal_2021} only considers binary treatments. In contrast, our method accommodates a finite number of treatments. 

% Paragraph 3: Limitations
Our proposed method has a few limitations. It assumes that one cause is more important than others, which could affect its applicability in cases where the relative importance of causes is unclear or varies. Additionally, our method does not include a tie-breaker mechanism, meaning optimal ITRs may not be unique. Finally, the user determines the pre-specified tolerance for our method, making it dependent on their judgment and expertise.
% Paragraph 4: Future direction
In this manuscript, we provide a meaningful advancement for decision-making in complex scenarios through multi-utility optimization. Future research may include incorporating tiebreakers for scenarios with multiple optimal ITRs or extending to multi-stage treatment settings. Such investigations would provide valuable insights into the method’s versatility. 
